
\chapter{Benchmark Environment}
\begin{spacing}{1.3}

All latency, throughput, and memory benchmarks were conducted on a controlled, reproducible hardware and software stack. Table~\ref{tab:bench_env} details the environment.

\begin{table}[H]
\centering
\caption{Benchmark environment specifications.}
\label{tab:bench_env}
\begin{tabular}{ll}
\toprule
\textbf{Component} & \textbf{Specification} \\
\midrule
\multicolumn{2}{c}{\textbf{Hardware -- Development (x86\_64)}} \\
\midrule
CPU                     & Intel Core i7-12700H (14 cores, 20 threads) \\
RAM                     & 32~GB DDR4-3200 \\
GPU (training)          & NVIDIA GeForce RTX 3060 (12~GB VRAM) \\
Storage                 & 1~TB NVMe SSD \\
\midrule
\multicolumn{2}{c}{\textbf{Hardware -- Edge Target (ARM64)}} \\
\midrule
Device                  & Raspberry Pi 4 Model B \\
SoC                     & Broadcom BCM2711, Cortex-A72 (4 cores at 1.8~GHz) \\
RAM                     & 8~GB LPDDR4 \\
Storage                 & 64~GB microSD (UHS-I) \\
\midrule
\multicolumn{2}{c}{\textbf{Software}} \\
\midrule
Operating system        & Ubuntu 22.04.3 LTS (x86\_64) / Raspberry Pi OS Bookworm (ARM64) \\
Python                  & 3.10.20 \\
PyTorch                 & 2.1.2+cu121 \\
ONNX Runtime            & 1.17.1 \\
ONNX                    & 1.15.0 \\
Execution provider      & CPUExecutionProvider (4 threads) \\
Docker                  & 24.0.7 (with Buildx + QEMU for ARM64) \\
\midrule
\multicolumn{2}{c}{\textbf{Benchmark Methodology}} \\
\midrule
Warm-up iterations      & 1,000 (discarded) \\
Measured iterations     & 10,000 \\
Input tensor            & float32, [1, 3, 224, 224] \\
Latency measurement     & Wall-clock, end-to-end including pre- and post-processing \\
Throughput calculation  & Batches per second times batch size, synchronous execution \\
Memory measurement      & Python tracemalloc + container-level via procfs \\
CPU governor            & Performance (frequency scaling disabled) \\
\bottomrule
\end{tabular}
\end{table}

All benchmarks were executed with no other significant user processes competing for CPU or memory resources. CPU frequency scaling was disabled via the performance governor to eliminate variability from dynamic frequency adjustment. The INT8 quantised model used static quantisation with a calibration dataset of 300 representative images drawn from the training set. Raspberry Pi~4 latency projections (15--18~ms p99) are based on the measured 3--4$\times$ performance ratio between x86\_64 and ARM Cortex-A72 CPUs at equivalent clock speeds for ONNX Runtime single-threaded inference.

\end{spacing}
